%% ============================================================
%%  CHAPTER 2 — RELATED WORK
%% ============================================================
\chapter{Related Work}
\label{ch:related_work}

This chapter reviews the existing body of research relevant to the
design and implementation of the proposed system. Section~\ref{sec:rw_bci_overview}
provides a broad overview of \acrshort{bci} paradigms and their
suitability for prosthetic control. Section~\ref{sec:rw_mi} examines
motor imagery-based \acrshort{bci} systems in depth.
Section~\ref{sec:rw_artifact} surveys artifact-based \acrshort{bci}
approaches — the paradigm adopted in this project.
Section~\ref{sec:rw_features} reviews feature extraction and
classification methods. Section~\ref{sec:rw_safety} discusses
safety mechanisms and discrete control strategies.
Section~\ref{sec:rw_3d} covers 3D-printed prosthetic technologies.
Section~\ref{sec:rw_gaps} synthesises the identified research gaps
and positions this project's contributions within the literature.

%% ----------------------------------------------------------
\section{Brain--Computer Interface Paradigms}
\label{sec:rw_bci_overview}
%% ----------------------------------------------------------

A \acrfull{bci} is a system that establishes a direct communication
channel between the brain and an external device by measuring,
processing, and translating neural activity into control
signals~\cite{wolpaw2002bci}. \acrshort{bci} systems are classified
along two primary dimensions: the invasiveness of the neural
recording method, and the type of neural signal exploited for
control.

\subsection{Recording Modalities}

Invasive \acrshort{bci} systems — including intracortical electrode
arrays and electrocorticography (\acrshort{eeg}) grids placed directly
on the cortical surface — offer the highest spatial resolution and
signal-to-noise ratio~\cite{hochberg2006neuronal, collinger2013high}.
Hochberg et al.~\cite{hochberg2006neuronal} demonstrated that a
tetraplegic patient could control a computer cursor and a robotic arm
using intracortical signals with remarkable precision. Collinger
et al.~\cite{collinger2013high} extended this to a 10-degree-of-freedom
robotic arm with coordinated, intuitive movements. However, invasive
approaches require neurosurgery, carry risks of infection and
electrode degradation, and demand lifelong clinical supervision —
making them impractical for the vast majority of potential users.

Non-invasive \acrshort{bci} modalities avoid surgical intervention
entirely. Functional near-infrared spectroscopy (fNIRS) measures
haemodynamic responses but suffers from a latency of several seconds,
making real-time control impractical~\cite{naseer2015fnirs}.
Magnetoencephalography (MEG) offers excellent temporal and spatial
resolution but requires magnetically shielded rooms and
superconducting sensors costing millions of dollars, precluding
portable deployment. \acrfull{eeg}, which measures scalp electrical
potentials generated by synchronised cortical neuron populations,
offers the most practical combination of temporal resolution
(millisecond-scale), cost (consumer devices available below
\$1{,}000), portability, and safety~\cite{lotte2007review}. These
properties make \acrshort{eeg} the dominant modality for assistive
\acrshort{bci} research and the natural choice for this project.

\subsection{Signal Paradigms}

\acrshort{eeg}-based \acrshort{bci} systems exploit one of several
neural signal paradigms:

\textbf{Steady-State Visual Evoked Potentials (SSVEPs)} are rhythmic
cortical responses elicited by flickering visual stimuli at specific
frequencies. The user attends to a stimulus flickering at a target
frequency, and the corresponding frequency component is detected in
the occipital \acrshort{eeg}. SSVEP systems achieve high information
transfer rates but require the user to maintain visual fixation on
a display, limiting their utility for users with impaired vision or
in mobile contexts~\cite{nakanishi2018ssvep}.

\textbf{P300 Event-Related Potentials} exploit a positive deflection
occurring approximately 300~ms after a rare, task-relevant stimulus.
P300-based \acrshort{bci}s can support large command vocabularies
but are slow (requiring multiple stimulus repetitions for reliable
detection) and stimulus-dependent~\cite{donchin2000p300}.

\textbf{Sensorimotor Rhythm / Motor Imagery} \acrshort{bci}s decode
voluntary modulations of mu (8--13~Hz) and beta (13--30~Hz)
oscillations generated during imagined or attempted
movements~\cite{pfurtscheller1999erd}. This paradigm supports
self-paced, asynchronous operation without external stimuli — a
critical advantage for prosthetic control. However, the subtle nature
of the underlying signals and high inter-subject variability impose
significant challenges, discussed in detail in
Section~\ref{sec:rw_mi}.

\textbf{Voluntary Artifact-Based Control} leverages large-amplitude
\acrshort{eeg} deflections produced by intentional facial and jaw
muscle activations. Although formally classified as artifacts rather
than neural signals, these signals are entirely suitable for
\acrshort{bci} control purposes when generated deliberately.
This paradigm is the focus of Section~\ref{sec:rw_artifact} and
the basis of the present work.

%% ----------------------------------------------------------
\section{Motor Imagery BCI Systems}
\label{sec:rw_mi}
%% ----------------------------------------------------------

Motor imagery (\acrshort{mi}) \acrshort{bci} is the most extensively
studied paradigm for prosthetic hand control. The neurophysiological
basis is well established: imagining a movement activates largely
the same cortical regions as executing it, producing detectable
Event-Related Desynchronisation (\acrshort{erd}) in the mu and beta
bands over the contralateral sensorimotor cortex and Event-Related
Synchronisation (\acrshort{ers}) over the ipsilateral
cortex~\cite{pfurtscheller1999erd}.

\subsection{Foundational Work}

Wolpaw et al.~\cite{wolpaw2002bci} established the theoretical
framework for motor imagery \acrshort{bci}, demonstrating that users
could learn to voluntarily modulate their sensorimotor rhythms for
cursor control. Pfurtscheller and Lopes da Silva~\cite{pfurtscheller1999erd}
characterised the \acrshort{erd}/\acrshort{ers} patterns in detail,
providing the neurophysiological foundation for spatial
discrimination between left- and right-hand imagery.
Blankertz et al.~\cite{blankertz2008csp} developed the Common Spatial
Patterns (\acrshort{csp}) algorithm into the de facto standard for
motor imagery feature extraction, demonstrating competition-winning
performance with a \acrshort{csp}--\acrshort{lda} pipeline on the
BCI Competition datasets.

\subsection{Prosthetic Control Applications}

M\"{u}ller-Putz et al.~\cite{mullerputz2008prosthetic} demonstrated
a four-class motor imagery system achieving 72--76\% accuracy for
prosthetic hand control, mapping imagined hand/wrist movements to
grasp types. Horki et al.~\cite{horki2011combined} achieved online
prosthetic grasping using combined motor imagery and SSVEP,
highlighting the value of hybrid approaches. Al-Timemy
et al.~\cite{al2013classification} investigated EEG-only control
as a fallback for amputees with insufficient residual EMG, confirming
feasibility at reduced accuracy compared to myoelectric approaches.

\subsection{Deep Learning Approaches}

Recent years have seen growing interest in applying deep learning
to motor imagery classification. Schirrmeister et al.~\cite{schirrmeister2017deep}
proposed deep convolutional neural networks (CNNs) that learn
spatial--temporal features directly from raw \acrshort{eeg} without
manual feature engineering, achieving 73--82\% on BCI Competition
benchmarks. Lawhern et al.~\cite{lawhern2018eegnet} developed EEGNet,
a compact CNN architecture with fewer than 2{,}000 parameters designed
specifically for \acrshort{eeg} \acrshort{bci} across multiple
paradigms. While these approaches show promise in offline benchmarks,
they require substantially larger training datasets (hundreds to
thousands of trials per class), GPU-level computation for training,
and have shown inconsistent advantages in online deployment scenarios
where data is limited and non-stationarity is present.

\subsection{Limitations of Motor Imagery BCI}

Despite its prominence, the motor imagery paradigm faces several
well-documented limitations that motivated the paradigm shift in
this project:

\begin{itemize}[leftmargin=*]
    \item \textbf{Signal amplitude.} \acrshort{mi}-related
    \acrshort{erd}/\acrshort{ers} modulations are typically
    1--5~\si{\micro\volt} in amplitude — comparable to the noise
    floor of consumer-grade \acrshort{eeg} devices — making reliable
    detection challenging~\cite{lotte2007review}.

    \item \textbf{BCI illiteracy.} Between 15\% and 30\% of
    individuals cannot reliably modulate their sensorimotor rhythms,
    regardless of training duration~\cite{vidaurre2010bci}. The
    neurological basis for this variability is not fully understood,
    but it represents a fundamental barrier to universal deployment.

    \item \textbf{Training burden.} Achieving reliable motor imagery
    control typically requires multiple calibration sessions totalling
    several hours, with performance degrading if sessions are
    separated by more than 24~hours due to \acrshort{eeg}
    non-stationarity~\cite{shenoy2006towards}.

    \item \textbf{Lab-to-real-world gap.} Vidaurre and
    Blankertz~\cite{vidaurre2010bci} documented systematic drops of
    10--20\% in classification accuracy between controlled laboratory
    recordings and real-world online operation, attributed to
    user fatigue, signal drift, and the absence of feedback during
    calibration.

    \item \textbf{Channel requirements.} Effective \acrshort{csp}
    feature extraction benefits from dense electrode coverage over
    the sensorimotor cortex, which consumer-grade devices (14--19
    channels) provide only partially, compared to the 64-channel
    research-grade systems on which most published benchmarks were
    obtained~\cite{blankertz2008csp}.
\end{itemize}

Table~\ref{tab:mi_comparison} summarises the performance of
representative motor imagery \acrshort{bci} systems from the
literature.

\begin{table}[htbp]
\centering
\caption{Comparison of motor imagery BCI systems from the literature.
All accuracy figures are offline unless stated otherwise.}
\label{tab:mi_comparison}
\renewcommand{\arraystretch}{1.3}
\begin{tabular}{p{2.8cm} c c p{2.8cm}}
\toprule
\textbf{Study} & \textbf{Classes} & \textbf{Accuracy} & \textbf{Method} \\
\midrule
Blankertz et al., 2008~\cite{blankertz2008csp}
    & 2--4 & 78--85\% & CSP + LDA \\
M\"{u}ller-Putz et al., 2008~\cite{mullerputz2008prosthetic}
    & 4 & 72--76\% & CSP + LDA \\
Horki et al., 2011~\cite{horki2011combined}
    & 4 & 70--75\% & CSP + Band Power + LDA \\
Schirrmeister et al., 2017~\cite{schirrmeister2017deep}
    & 4 & 73--82\% & Deep CNN \\
Lawhern et al., 2018~\cite{lawhern2018eegnet}
    & 4 & 75--82\% & EEGNet (compact CNN) \\
\textbf{This work}
    & \textbf{5} & \textbf{91.3\%} & \textbf{Hybrid (features + DTW + Riemann)} \\
\bottomrule
\end{tabular}
\end{table}

%% ----------------------------------------------------------
\section{Artifact-Based BCI Systems}
\label{sec:rw_artifact}
%% ----------------------------------------------------------

Voluntary physiological artifacts — \acrshort{eeg} signals generated
by intentional facial, ocular, and jaw muscle contractions — have
received considerably less attention than motor imagery in the
\acrshort{bci} literature, despite offering several practical
advantages. The growing body of work in this area can be divided
into three categories: eye blink-based control, jaw/facial EMG
control, and combined multi-artifact systems.

\subsection{Eye Blink-Based Control}

Eye blinks produce characteristic large-amplitude deflections
(50--200~\si{\micro\volt}) predominantly at frontal electrodes
(Fp1, Fp2) due to the corneoretinal potential and associated
frontalis muscle activity~\cite{usakli2010artifact}. The blink
waveform is dominated by a sharp positive peak followed by a slower
return to baseline, with the entire event lasting approximately
200--400~ms.

Usakli et al.~\cite{usakli2010artifact} provided one of the first
systematic studies of voluntary eye blink as a \acrshort{bci}
control signal, demonstrating that single and double blinks could
be reliably distinguished using amplitude and temporal features
extracted from frontal electrodes, with classification accuracies
exceeding 90\% in offline analysis. Zander et al.~\cite{zander2011bci}
explored the use of blinks as a ``binary switch'' for \acrshort{bci}
applications requiring only an on/off command, demonstrating
robustness to electrode impedance variations.

Distinguishing single from double blinks — as required in this
project — presents a significantly harder classification problem
than simple blink detection. The inter-blink interval, the number
of deflection peaks in the frontal signal, and the total event
duration are the natural discriminating features. However, as
this project discovered, event-counting approaches are fragile when
the recorded single-blink data is contaminated with involuntary
multi-blink clusters; a shape-based approach using Dynamic Time
Warping (\acrshort{dtw}), discussed in Section~\ref{subsec:dtw_litt},
proves substantially more robust. To the best of the authors'
knowledge, no published work has applied \acrshort{dtw}
shape-matching to voluntary blink-count discrimination as a
\acrshort{bci} command set, representing a novel contribution of
this project.

\subsection{Jaw Clinch and Bruxism Control}

Jaw clinching (sustained bilateral biting force) and bruxism
(lateral grinding movements) produce strong \acrshort{emg} signals
in the masseter and temporalis muscles, which contaminate the
\acrshort{eeg} primarily at temporal electrodes (T3, T4) in the
international 10--20 system~\cite{croft2000removal}. The
\acrshort{emg} contamination is broadband (30--300~Hz) with
significant power extending into the \acrshort{eeg} frequency range,
producing large-amplitude artifacts at temporal sites.

Fatourechi et al.~\cite{fatourechi2007emg} characterised jaw
\acrshort{emg} artifacts in \acrshort{eeg} recordings and proposed
methods for their removal — the opposite of the use case in this
project, where the artifact is the desired signal. Kilicarslan
et al.~\cite{kilicarslan2016high} demonstrated a wearable
\acrshort{bci} using jaw clenching as one of several control
modalities, achieving reliable command detection in ambulatory
conditions.

The discrimination of left versus right bruxism — a key challenge
in this project — exploits the spatial lateralisation of the
temporalis muscle activation. Right-sided bruxism produces stronger
\acrshort{emg} contamination at T4 (right temporal), while
left-sided bruxism dominates at T3 (left temporal). This asymmetry
provides the discriminating feature but is confounded by the
bilateral nature of some jaw movements and the proximity of T3
and T4 to the clinch artifact, which activates both temporal
sites simultaneously.

\subsection{Multi-Artifact Command Systems}

Hybrid systems combining multiple voluntary artifact types to create
larger command vocabularies have been explored in the context of
alternative communication aids for severely paralysed users.
Huang et al.~\cite{huang2019hybrid} combined eye blinks, jaw
clenching, and head movements into a six-command communication
interface, demonstrating the feasibility of multi-artifact command
sets. However, their work focused on communication rather than
prosthetic control, and did not address the specific challenges
of real-time prosthetic command execution, including the activity
gate problem, cross-session normalisation, or FSM-based safety
validation.

Pfurtscheller et al.~\cite{pfurtscheller2003thought} demonstrated
that a combination of \acrshort{eeg} and \acrshort{emg} signals
could restore hand grasp function in a tetraplegic patient via
functional electrical stimulation (FES), establishing a precedent
for using non-motor-imagery neural and para-neural signals for
prosthetic actuation.

Table~\ref{tab:artifact_comparison} compares selected artifact-based
\acrshort{bci} systems from the literature.

\begin{table}[htbp]
\centering
\caption{Comparison of artifact-based BCI systems from the
literature.}
\label{tab:artifact_comparison}
\renewcommand{\arraystretch}{1.3}
\begin{tabular}{p{2.5cm} p{2.2cm} c p{2.5cm}}
\toprule
\textbf{Study} & \textbf{Artifact Type} & \textbf{Commands} &
\textbf{Accuracy} \\
\midrule
Usakli et al., 2010~\cite{usakli2010artifact}
    & Eye blink & 2 & $>$90\% \\
Zander et al., 2011~\cite{zander2011bci}
    & Eye blink & 2 (switch) & $>$85\% \\
Kilicarslan et al., 2016~\cite{kilicarslan2016high}
    & Jaw clench & 3 & $\sim$88\% \\
Huang et al., 2019~\cite{huang2019hybrid}
    & Blink + jaw + head & 6 & $\sim$82\% \\
\textbf{This work}
    & \textbf{Blink (2) + jaw + bruxism (2)}
    & \textbf{5} & \textbf{91.3\%} \\
\bottomrule
\end{tabular}
\end{table}

%% ----------------------------------------------------------
\section{Feature Extraction and Classification Methods}
\label{sec:rw_features}
%% ----------------------------------------------------------

\subsection{Time-Domain Features}

Time-domain features characterise the statistical and morphological
properties of \acrshort{eeg} epochs directly from the raw signal
samples. Commonly used time-domain features include the Root Mean
Square (\acrshort{rms}) amplitude, Peak-to-Peak (\acrshort{ptp})
amplitude, zero-crossing rate, and Hjorth
parameters~\cite{hjorth1970eeg}.

Hjorth parameters — activity, mobility, and complexity — provide
a compact and computationally efficient characterisation of signal
dynamics. Activity measures the signal variance (power), mobility
captures the mean frequency, and complexity describes the bandwidth
of the signal relative to a pure sine wave. These parameters were
originally proposed by Hjorth~\cite{hjorth1970eeg} and have since
been widely used in \acrshort{eeg} analysis for their robustness
to noise and their lack of dependence on spectral estimation
parameters.

\subsection{Frequency-Domain Features}

Band power features quantify the energy within specific frequency
bands using Power Spectral Density (\acrshort{psd}) estimation.
For motor imagery \acrshort{bci}, the standard bands are mu
(8--13~Hz) and beta (13--30~Hz). For artifact-based \acrshort{bci},
the relevant bands extend to the gamma range ($>$30~Hz) and beyond,
since jaw \acrshort{emg} artifacts contain substantial energy in
the 30--80~Hz range. Relative band power — computed as the ratio
of band power to total power — provides amplitude-invariant features
that generalise better across sessions with varying electrode
impedance~\cite{lotte2007review}.

\subsection{Common Spatial Patterns}

\acrfull{csp} is the standard spatial filtering technique for
motor imagery \acrshort{bci}~\cite{blankertz2008csp}. \acrshort{csp}
solves a generalised eigenvalue problem to find spatial filters
(linear combinations of channels) that simultaneously maximise
variance for one class while minimising it for the other. The
resulting log-variance features are highly discriminative for
two-class motor imagery problems. For multi-class problems,
pairwise \acrshort{csp} or Filter Bank \acrshort{csp}
(FBCSP)~\cite{ang2012fbcsp} extensions are typically employed.

For the artifact-based paradigm in this project, \acrshort{csp}
is less directly applicable because the discriminating information
is localised to specific electrode groups (frontal for blinks,
temporal for jaw) rather than distributed across the cortex in
a pattern that \acrshort{csp} is designed to extract. Instead, a
channel-aware feature engineering approach — selecting features
from physiologically motivated electrode groups — was adopted,
as detailed in Chapter~6.

\subsection{Classifier Comparison}

\textbf{Linear Discriminant Analysis (\acrshort{lda})} is the most
widely used classifier for motor imagery \acrshort{bci} due to its
computational simplicity, interpretability, and effectiveness with
\acrshort{csp} features and limited training
data~\cite{lotte2007review}. \acrshort{lda} assumes Gaussian class
distributions with equal covariance matrices, an assumption that
holds reasonably well for \acrshort{csp} log-variance features but
may be violated for the heterogeneous artifact feature sets used
in this project.

\textbf{Support Vector Machines (\acrshort{svm})} with nonlinear
kernels can outperform \acrshort{lda} when class boundaries are
non-linear, at the cost of increased computational complexity and
hyperparameter sensitivity~\cite{lal2004svm}. For real-time
applications with limited training data, the performance advantage
of \acrshort{svm} over \acrshort{lda} is typically marginal.

\textbf{Random Forests (\acrshort{rf})} are ensemble methods that
combine multiple decision trees trained on random feature subsets
and bootstrap samples~\cite{breiman2001rf}. Random Forests are
robust to irrelevant features, handle mixed feature types naturally,
provide feature importance estimates, and are largely insensitive
to hyperparameter choices. These properties make them well suited
to the heterogeneous 93-feature vector in this project, which
combines time-domain, frequency-domain, and asymmetry features of
varying scales and distributions.

Table~\ref{tab:classifier_comparison} summarises the properties of
the three classifier families considered in this project.

\begin{table}[htbp]
\centering
\caption{Comparison of classifier properties relevant to this
project.}
\label{tab:classifier_comparison}
\renewcommand{\arraystretch}{1.3}
\begin{tabular}{p{2.0cm} p{2.2cm} p{1.5cm} p{2.5cm}}
\toprule
\textbf{Classifier} & \textbf{Boundary} &
\textbf{Training Data} & \textbf{Key Properties} \\
\midrule
LDA & Linear & Low & Fast, interpretable, assumes Gaussian \\
SVC (RBF) & Non-linear & Medium & Better boundaries, needs tuning \\
Random Forest & Non-linear & Medium & Robust, feature importance,
    handles mixed features \\
\bottomrule
\end{tabular}
\end{table}

\subsection{Dynamic Time Warping for Waveform Classification}
\label{subsec:dtw_litt}

Dynamic Time Warping (\acrshort{dtw}) is a classical algorithm for
measuring the similarity between two temporal sequences that may vary
in speed or timing~\cite{sakoe1978dtw}. Rather than comparing samples
at fixed time indices, \acrshort{dtw} finds an optimal elastic
alignment that warps the time axis of one sequence onto the other,
minimising the cumulative distance between aligned points. This makes
it well suited to comparing biological waveforms whose morphology is
consistent but whose exact timing varies from instance to instance.

In the \acrshort{bci} and biosignal literature, \acrshort{dtw} has
been applied to gesture and gait recognition, where the \emph{shape}
of a movement signal is more reliable than its precise timing or
amplitude~\cite{berndt1994dtw}. When combined with a
$k$-nearest-neighbour (\acrshort{knn}) classifier --- assigning each
test sequence the majority label of its closest exemplars under the
\acrshort{dtw} distance --- the approach requires no parametric model
and is robust to noisy or partially contaminated training data, since
classification depends on proximity to individual exemplars rather
than to a class average. This project applies \acrshort{dtw}--\acrshort{knn}
to blink-count discrimination, where exactly this robustness to
contaminated single-blink recordings proves decisive (Chapter~9).

\subsection{Riemannian Geometry for EEG Classification}
\label{subsec:riemann_litt}

A more recent and increasingly influential family of \acrshort{bci}
classification methods operates on the \emph{covariance matrices} of
multichannel signals rather than on hand-crafted features. Because the
spatial covariance between channels captures how signal energy is
distributed across the scalp, it encodes discriminative spatial
structure while being inherently invariant to certain global amplitude
changes. Crucially, covariance matrices are symmetric positive-definite
(SPD) and therefore live on a curved Riemannian manifold rather than in
flat Euclidean space; treating them with the correct Riemannian
geometry, rather than naively as flat vectors, yields substantially
better classification~\cite{barachant2012riemann}.

The standard pipeline estimates a covariance matrix per epoch, then
either classifies directly by Minimum Distance to Mean
(\acrshort{mdm}) under the affine-invariant Riemannian metric, or
projects the matrices into a Euclidean \emph{tangent space} where a
conventional classifier such as an \acrshort{svc} can be
applied~\cite{barachant2012riemann}. Riemannian methods have won
multiple \acrshort{bci} competitions and are widely regarded as
state-of-the-art for motor imagery. Their key practical advantage for
this project is robustness to cross-session amplitude drift: because
the discriminative information lives in the \emph{structure} of the
covariance rather than its absolute scale, Riemannian classifiers
generalise across recording days far better than amplitude-based
features. This property is exploited for the muscle-command branch of
the hybrid architecture (Chapter~9).

\subsection{Class Imbalance Handling}

Class imbalance is a common challenge in \acrshort{bci} datasets,
where some artifact types are easier or more natural to perform
repeatedly than others, leading to unequal epoch counts per class.
Training a classifier on imbalanced data produces models biased
toward majority classes, inflating overall accuracy while degrading
performance on minority classes. The \acrfull{smote}~\cite{chawla2002smote}
addresses this by generating synthetic minority class samples through
interpolation in feature space between existing minority class
examples, effectively balancing the training set without discarding
majority class data.

%% ----------------------------------------------------------
\section{Safety Mechanisms and Discrete Control}
\label{sec:rw_safety}
%% ----------------------------------------------------------

The safety of \acrshort{bci}-controlled prosthetic devices is a
critical concern that is often inadequately addressed in academic
literature. Direct execution of classifier outputs without any
validation layer leads to unstable, erratic prosthetic behaviour,
as transient misclassifications — which occur even in high-accuracy
systems — translate directly into unintended physical movements.

\subsection{Temporal Confirmation}

Gao et al.~\cite{gao2021interface} proposed temporal confirmation
as a general mechanism for reducing false positive command execution
in \acrshort{bci} systems. By requiring the classifier to produce
the same prediction across multiple consecutive decision windows
before executing the corresponding command, transient
misclassifications are filtered out. Their analysis showed that
requiring three consecutive matching predictions reduced false
positive rates from 15--20\% to below 5\% at the cost of a moderate
increase in response latency.

\subsection{Finite State Machine Control}

Mil\'{a}n et al.~\cite{milan2010combining} proposed Finite State
Machine (\acrshort{fsm}) architectures for \acrshort{bci} prosthetic
control, encoding the logical constraints on valid command sequences
as state transition rules. An \acrshort{fsm} ensures that the
prosthetic hand can only transition between states in physically
meaningful ways — for example, preventing a direct transition from
``rotate left'' to ``rotate right'' without passing through a neutral
state, which would cause mechanical shock and potentially damage
the actuators.

\subsection{Discrete vs.\ Continuous Control}

M\"{u}ller-Putz and Pfurtscheller~\cite{mullerputz2008prosthetic}
demonstrated empirically that discrete command-based prosthetic
control achieves more stable real-world performance than continuous
proportional control with non-invasive \acrshort{eeg}, for two
related reasons: first, the limited information transfer rate of
non-invasive \acrshort{eeg} is better matched to a small set of
discrete states than to continuous multi-dimensional trajectories;
second, discrete commands are robust to brief classification errors,
whereas continuous control amplifies any noise in the classifier
output into visible, unintended movements.

%% ----------------------------------------------------------
\section{3D-Printed Prosthetic Technologies}
\label{sec:rw_3d}
%% ----------------------------------------------------------

The democratisation of prosthetic technology through additive
manufacturing (3D printing) has been a significant development over
the past decade. Open-source designs such as e-NABLE~\cite{enable2015}
and InMoov~\cite{langevin2026inmoov} provide freely available, fully
documented prosthetic hand designs printable on standard desktop
FDM printers using PLA filament, reducing fabrication costs from
\$10{,}000--\$100{,}000 for commercial myoelectric hands to
\$100--\$500.

Ten Kate et al.~\cite{ten2017review} conducted a systematic review
of 3D-printed upper-limb prostheses, finding that while FDM-printed
PLA devices are unsuitable for heavy-duty functional tasks due to
anisotropic mechanical properties and layer adhesion limitations,
they are entirely adequate for research prototypes and technology
demonstration — the role they serve in this project.

The InMoov design, upon which the prosthetic hand in this project
is based, provides individual finger actuation through a
tendon-driven mechanism using monofilament fishing line routed
through printed pulleys to servo motors housed in the forearm.
The modular architecture facilitates assembly, maintenance, and
customisation. Integration of \acrshort{bci} control with
InMoov-based designs has been demonstrated in several university
research projects~\cite{langevin2026inmoov}, though published
documentation of the specific control interface, Arduino firmware,
and real-time communication protocol is sparse.

%% ----------------------------------------------------------
\section{Research Gaps and Project Positioning}
\label{sec:rw_gaps}
%% ----------------------------------------------------------

The literature review reveals the following critical gaps that
directly motivate the design decisions in this project:

\begin{enumerate}[leftmargin=*, label=\textbf{G\arabic*.}]

    \item \textbf{Physiology-matched hybrid architectures are
    unexplored.} Existing artifact-based \acrshort{bci} systems apply
    a single flat classifier to all gesture classes, despite those
    classes arising from fundamentally different signal types
    (ocular \acrshort{eog} for blinks, muscular \acrshort{emg} for
    jaw gestures). No published work routes each gesture family to a
    classification method matched to its physiology --- for example,
    shape-matching for blink waveforms and covariance geometry for
    muscle activations --- within a single integrated system.

    \item \textbf{Resting window contamination is not addressed.}
    None of the reviewed artifact-based \acrshort{bci} studies
    explicitly address the problem of rest windows within gesture
    recordings carrying the gesture class label. This contamination
    — which we show reduces accuracy by 3.1\% — has been ignored
    in the literature, likely because most studies use marker-aligned
    epoching rather than continuous sliding window extraction.

    \item \textbf{Data leakage in overlapping epoch evaluation is
    underreported.} The data leakage introduced by splitting
    overlapping epochs randomly into training and test sets is
    acknowledged in the signal processing literature but rarely
    discussed explicitly in \acrshort{bci} papers, leading to
    potentially inflated accuracy reports across the field.

    \item \textbf{Lateral bruxism asymmetry features are novel.}
    While jaw \acrshort{emg} artifacts have been characterised in
    the literature, the use of a normalised T3--T4 asymmetry index
    as an explicit feature for left/right bruxism discrimination
    has not been reported. Existing work either treats bruxism as
    a single undifferentiated artifact class or uses raw channel
    amplitudes without the normalised asymmetry formulation.

    \item \textbf{End-to-end artifact BCI with FSM and activity
    gate is missing.} No published system combines: (a) a
    multi-class voluntary artifact command set, (b) activity-gated
    training epoch extraction, (c) a real-time sliding window
    inference engine, and (d) an \acrshort{fsm} safety layer with
    confidence thresholding — in a single integrated prosthetic
    control system.

    \item \textbf{Consumer hardware validation is limited.} The
    majority of artifact-based \acrshort{bci} studies use
    research-grade 32--64 channel \acrshort{eeg} amplifiers.
    Validation on consumer-grade hardware with fewer channels
    and lower signal quality — the hardware most accessible to
    users in resource-constrained settings — is rare.

\end{enumerate}

This project directly addresses gaps G1 through G5, and partially
addresses G6 by using a 19-channel device representative of
mid-range consumer \acrshort{eeg} hardware. Table~\ref{tab:gap_positioning}
summarises how the contributions of this work map onto the
identified gaps.

\begin{table}[htbp]
\centering
\caption{Mapping of project contributions to identified research
gaps.}
\label{tab:gap_positioning}
\renewcommand{\arraystretch}{1.3}
\begin{tabular}{c p{4.5cm} p{4.0cm}}
\toprule
\textbf{Gap} & \textbf{Description} & \textbf{Contribution} \\
\midrule
G1 & Physiology-matched hybrid architecture & Hybrid routing: features + DTW + Riemann (C1) \\
G2 & Rest window contamination & Activity-gated epoching (C4) \\
G3 & Data leakage in evaluation & Leakage audit + LOSO (C5) \\
G4 & L/R bruxism asymmetry feature & T3--T4 asymmetry index (C2) \\
G5 & End-to-end artifact BCI + FSM & Full integrated system (C6) \\
G6 & Consumer hardware validation & 19-ch device, 200~Hz (C7) \\
\bottomrule
\end{tabular}
\end{table}

\section{Chapter Summary}
\label{sec:rw_summary}

This chapter reviewed the relevant literature across six areas:
\acrshort{bci} paradigms, motor imagery systems, artifact-based
control, feature extraction and classification methods, safety
mechanisms, and 3D-printed prosthetics. The motor imagery paradigm,
while dominant, faces fundamental limitations in signal amplitude,
user variability, and training burden that limit its practical
deployment. Artifact-based \acrshort{bci} offers larger, more
reproducible signals but remains underexplored, particularly for
multi-class command sets and real-time prosthetic control. Six
specific research gaps were identified, all of which are addressed
by the contributions of this project. The following chapters present
the complete system design, implementation, and validation in detail.